\section{The actual PEL uniformization at an arithmetic point}
\label{sec:dvr-pel-uniformization}

The individual integral model now gives the required degeneration at
every arithmetic point.  We verify its period, line, section and
torsion interfaces.  The original level even supplies a split torus
over the same discrete valuation ring.

Let $R$ be a complete discrete valuation ring of characteristic zero,
with fraction field $K$, residue field $\kappa$ of characteristic $p$,
and uniformizer $\pi$.  Let $c.F\hookrightarrow K$ be the original
coefficient embedding, and let $\mathcal A/K$ be an original moduli
object at a level satisfying $\mathrm{GoodAt}(p)$.
Pull back the individual projective model of
Theorem~\ref{thm:original-integral-pel-model} along the unique
valuative extension of this point.  Denote its semiabelian family by
$\mathcal G/R$ and its original primitive cubical line by $\mathscr L$.
The same argument applies to any semiabelian model of the original
object: homomorphisms of semiabelian schemes over a normal base extend
uniquely from a dense open, and the original cubical lines have unique
extensions\cite[Propositions 3.3.1.5 and 3.3.2.2]{dvr-unif-lan}.

\begin{theorem}[Actual arithmetic-point supplier]
\label{thm:dvr-original-pel-unif}
Over this same $R$, the original model has a
\texttt{PELUnif} with \texttt{TorsionIso}.  Its line is precisely
$\mathscr L$, with the given zero rigidification.  Its Fourier
expansions reconstruct every integral thickening and retain one
common denominator for every generic section.  The torsion
identification holds for every integer $k\geq1$, including multiples
of the residue characteristic, and over every algebraically closed
extension of $K$.

Moreover $H^0(\mathcal G,\mathscr L)$ is a free rank-one $R$-module
whose generator is nonzero on the closed fibre.  Thus the original
\texttt{PushInvertible} condition holds over this discrete valuation
base.  Theorem~\ref{thm:pel-integral-theta} supplies the corresponding
statement over the entire higher-dimensional model.
\end{theorem}

\subsection*{The level splits the torus without enlarging the field}

Choose an odd prime $q\mid M$ with $q\ne p$, which exists by
$M/p^{v_p(M)}\geq3$.  The functorial half-to-full equivalence of
Theorem~\ref{thm:original-half-level-completion}, after the given
embedding $c.F\hookrightarrow K$, makes $\mathcal A[q]$ a constant
group over $K$.  The prescribed $\zeta_q\in c.F$ is also in $R$:
it and its inverse are integral over $R$.  Its reduction remains
primitive since $q$ is a unit.  Therefore $\mu_q$ is constant over
$R$ as well.

The kernel $\mathcal G[q]$ is quasi-finite and flat over $R$.
Its torus part $\mathcal G[q]^\mu$ is a finite flat subgroup over
the \emph{whole} henselian base, with closed fibre $T_\kappa[q]$;
it is canonically the $q$-torsion of the Raynaud torus
\cite[Section 3.4.1 and (3.4.2.1)]{dvr-unif-lan}.
Since $q$ is invertible, it is finite \'etale.  Its generic fibre is
a subgroup of the constant $\mathcal A[q]$, so is constant.  A finite
\'etale scheme over a normal discrete valuation ring with constant
generic fibre is constant: each generic component has function field
$K$, and its finite normal closure is $\operatorname{Spec}R$.
Thus $T_\kappa[q]$ is constant.

For the character lattice $X$ of $T_\kappa$, the constancy of both
$T_\kappa[q]$ and $\mu_q$ says that the Galois action on $X/qX$ is
trivial.  Its action on $X$ has finite image.  The principal
$q$-congruence subgroup of $\operatorname{GL}(X)$ has no torsion:
for odd $q$, the $q$-adic logarithm is injective on
$1+q\operatorname{End}(X\otimes\mathbb Z_q)$, and a finite-order
element has logarithm zero.  Hence the entire Galois action is
trivial.  The closed torus and the Raynaud torus are split over the
original base.  This includes $X=0$.

\subsection*{The exact principal line and period data}

The actual cubical line $\mathscr L$ is ample over $R$.
Indeed it is symmetric, its generic restriction is ample, and for
a symmetric cubical line on a semiabelian scheme over the normal
base, generic ampleness is equivalent to ampleness
\cite[Lemma 4.2.1.6]{dvr-unif-lan}.
The closed torus has just been shown split, so the associated formal
line algebraizes on the Raynaud extension by
\cite[Proposition 3.3.3.11]{dvr-unif-lan}.
Thus the hypotheses of the \emph{ample} degeneration equivalence
are met for this same line, before we apply that equivalence.

Raynaud's construction and the equivalence for polarized degenerations
give
\[
 0\longrightarrow T\longrightarrow\mathcal G^\natural
   \xrightarrow{\,q_B\,}\mathcal B\longrightarrow0,
 \qquad \iota:Y\longrightarrow\mathcal G^\natural_K,
 \qquad \phi:Y\longrightarrow X,
\]
with $X,Y$ free lattices of finite rank and $\mathcal B/R$ abelian
\cite[Sections 3.3.3, 4.2.1 and Theorem 4.4.16]{dvr-unif-lan}.
If $\lambda$ and $\lambda_B$ are the generic and abelian-part
polarizations, the exact degree formula is
\[
       \deg\lambda=[X:\phi(Y)]^2\deg\lambda_B.
\]
The square is essential\cite[Lemma 3.4.4.3]{dvr-unif-lan}.
Since the original polarization is principal, both positive integers
on the right are one.  We therefore identify $Y$ with $X$ using
\emph{this} $\phi$, and $\lambda_B$ is principal.  The period lattice
is consequently split too.  Its rank is retained throughout.

Write $w=w_0$ for the original element with $w+\bar w=1$.
The line on the abelian part is
\[
  \mathscr M=(\operatorname{id},\lambda_B\iota_B(w))^*
                    \mathcal P_B.
\]
It is the same primitive, principally polarizing line as in
Proposition~\ref{prop:pel-primitive-line}.  The cubical line
$\mathscr L^\natural$ on the Raynaud extension has
\[
 \delta\mathscr L^\natural
   =(q_B\times q_B)^*(\operatorname{id}\times\lambda_B)^*\mathcal P_B.
\]
Transport the \emph{original} primitive identity through the
equivalence of cubical degeneration data and pull back along the
graph of $w$.  This gives
$\mathscr L^\natural\simeq q_B^*\mathscr M$ with its original
rigidification and period action.  This is not a choice of an
arbitrary trivialization on the torus: its possible character
ambiguity is fixed by the transported primitive identity.
The equivalence respects tensor operations and homomorphisms
\cite[Theorem 4.4.16 and Section 5.1.1]{dvr-unif-lan}.

Let $c:X\to\mathcal B(R)$ be the resulting period point map.
The period factors are the original generators
\[
 \psi_y\in\mathscr M_{c(y)}^{-1}\otimes_R K,
 \qquad
 \tau_{y,\chi}
   =\psi_{y+\chi}\psi_y^{-1}\psi_\chi^{-1}\mathrm{rig}^{-1},
\]
with the indicated line contractions.  They are even and bimultiplicative
in precisely the source's cubical sense.  In particular the
coefficient-ring action is the original one: $c(ay)=\iota_B(a)c(y)$,
and its action on torus characters is $\bar a$, as follows from
Rosati adjunction.  This gives every action identity in
\texttt{DegPre.OAct} and \texttt{UnifPre.OLinear}.

For clarity, the transported primitive identity gives an integral
generator $u$ of the line on $\mathcal B$ with pullback maps
$(1+w,1,w,0)$ and exponents $(1,-2,-1,2)$, normalized at zero, and
the exact relation
\[
       \tau_{y,wy}\,u(c(y))=\psi_y\,\mathrm{rig}.
\]
This is the source's \texttt{TraceFormTwistIso}; the equality is an
identity of period descent data, with no additional character or
scalar.  It supplies its \texttt{w0} field.

Choose integral generators of the fibre lines when computing
valuations.  Symmetry and the cube relation imply that
$Q(y)=v(\psi_y)$ is an integer-valued quadratic form, with
$b(y,\chi)=v(\tau_{y,\chi})$ its polarization.  The Riemann
positivity condition gives $2Q(y)>0$ for $y\ne0$; hence $Q$ is
positive definite.  It follows that $\psi_y\in\pi R$ for nonzero
$y$ and, for every $m$, all but finitely many $\psi_y$ belong to
$\pi^mR$.  These are the exact \texttt{psi\_pos} and
\texttt{psi\_tendsto} fields, not a uniform lower bound independent
of the arithmetic point
\cite[Definitions 4.2.1.10--11 and Theorem 4.2.1.14]{dvr-unif-lan}.
The valuation is always $v(\pi)=1$; no integral division by two is used.

\subsection*{Integral expansions before the dimension argument}

On each thickening $R/\pi^{e+1}$, the formal Raynaud isomorphism
identifies $\mathcal G$ with the corresponding torus extension of
$\mathcal B$ and identifies $\mathscr L$ with $q_B^*\mathscr M$.
The character-$\chi$ coefficient line at level $k$ is exactly
\[
 \mathscr M_{k,\chi}
 =\mathscr M^{k-1}\otimes t_{c(\chi)}^*\mathscr M
        \otimes\mathscr M_{c(\chi)}^{-1}\otimes\mathscr M_0.
\]
Thus the four exponents remain $(k-1,1,-1,1)$.  The coordinate
algebra of the split torus torsor gives a direct sum over characters
on every thickening.  Its transitions and its normalization at zero
give the source's \texttt{GtCoord} and the compatible integral line
isomorphisms, including their symmetry and cubical identities.

Now start with an actual generic section
$s\in H^0(\mathcal A,\mathscr L_K^k)$.  Quasi-compactness of
$\mathcal G$ supplies \emph{one} $r\in R\setminus\{0\}$ such that
$r s$ extends to an integral section $\sigma$ on $\mathcal G$.
Expand $\sigma$ on all thickenings.  For each $\chi$, its compatible
coefficients are an element of
$H^0(\mathcal B,\mathscr M_{k,\chi})$ by proper formal functions on
$\mathcal B$.  On each thickening only finitely many coefficients
are nonzero, and their finite sum reconstructs $\sigma$ there.
After dividing by this same $r$, one obtains the generic coefficient
family of $s$.  In particular this is exactly the source's
\texttt{Expands}: a single denominator, an actual integral section,
integral coefficient sections, and reconstruction on every
$R/\pi^{e+1}$.  No integral conclusion is inferred from equal generic
dimensions\cite[Section 4.3.1]{dvr-unif-lan}.

In the original line-contraction notation the period recurrence is
\[
 s_{\chi+ky}
 =\bigl(\psi_y^k\tau_{y,\chi}\mathrm{rig}^{\,k}\bigr)
                       t_{c(y)}^*s_\chi.
\]
It is the defining \texttt{FamFE}, with the same signs and powers.
The Fourier map is injective by the formal Raynaud comparison.
Let $t=\operatorname{rank}X$ and $b=\dim\mathcal B$, so
$g=\dim\mathcal A=t+b$.  A family satisfying this recurrence is
uniquely and freely determined by its coefficients at representatives
of $X/kX$: consistency of successive transports is exactly the
cubical period cocycle.  Each coefficient space has dimension $k^b$,
because its line is $\mathscr M^k$ tensored with a line in
$\operatorname{Pic}^0(\mathcal B_K)$, with principal polarization.
Consequently the space of these families has dimension
\[
                  |X/kX|\,k^b=k^{t+b}=k^g.
\]
The generic section space also has dimension $k^g$.  The injective
Fourier map is therefore onto.  Since its expansions were constructed
with the integral common-denominator data first, this proves the
source's \texttt{ExpBij} for every $k\geq1$.  The same construction
commutes with field extension: generic sections and their coefficient
sections satisfy flat base change on the proper abelian varieties.
In particular it supplies the original \texttt{ExpandsF}, retaining
the specified line comparison rather than changing its scalar.

\subsection*{Torsion and its exact theta action, including $p$-torsion}

The source's \texttt{GtPtF} is the point description of
$\mathcal G^\natural$: over $b\in\mathcal B$, its coordinates are
generators of the inverse biextension fibres
$\delta\mathscr M(c(\chi),b)^{-1}$.  Their exponents
$(-1,1,1,-1)$ and multiplication are the original cube law.
This agrees with the minus-pushout convention for the Raynaud
extension\cite[Definition 4.2.1.13]{dvr-unif-lan}.

Accordingly a source \texttt{QkPt} is a pair $(y,z)$ satisfying
$[k]z=\iota(y)$, and its equivalence relation is
\[
              (y,z)\sim(y+kx,z+\iota(x)),\qquad x\in X.
\]
Mumford's torsion construction gives an isomorphism from this quotient
to $\mathcal A[k]$.  The construction is valid for every $k\geq1$:
\cite[Theorem 4.5.3.10]{dvr-unif-lan} describes the actual quasi-finite
group scheme $\mathcal G[k]$ using the period division pieces.
It does not require $k$ to be invertible in $R$.  We use its generic
fibre, and do not assert that all of $\mathcal G[k]$ is finite over
$R$.  Addition and the coefficient action commute with the quotient.
All the division points involved are algebraic over $K$, so the same
description holds over every algebraically closed extension field.

It remains to verify the source's stronger definition of the
identification by theta action.  Put
$\mathscr P(u,x)=\delta\mathscr M(u,x)$, and let $b=q_B(z)$,
so $kb=c(y)$.  In this notation the coefficient line is
$\mathscr C_{k,\alpha}(x)=\mathscr M_x^k\otimes
\mathscr P(c(\alpha),x)$.  The cube identities give the following
\emph{constant} ratio line on $\mathcal B_F$:
\[
 \begin{split}
 \mathscr C_{k,\alpha+y}\otimes
           (t_b^*\mathscr C_{k,\alpha})^{-1}
 &\simeq
 \underline{\mathscr P(c(\alpha),b)^{-1}
                  \otimes\mathscr M_0^k\otimes\mathscr M_b^{-k}}.
 \end{split}
\]
Indeed $\mathscr P(c(y),x)=\mathscr P(b,x)^k$ and
$\mathscr P(c(\alpha),x+b)=
\mathscr P(c(\alpha),x)\mathscr P(c(\alpha),b)$;
substituting
$\mathscr M_{x+b}=\mathscr M_x\mathscr M_b\mathscr M_0^{-1}
\mathscr P(b,x)$ cancels every $x$-dependent factor.

The coordinate $z.s(\alpha)$ is the original generator of
$\mathscr P(c(\alpha),b)^{-1}$.  Evaluating a theta lift at zero
provides one generator
$w\in\mathscr M_0^k\otimes\mathscr M_b^{-k}$.
Let $u_\alpha$ be the ratio-line generator whose zero value is
$z.s(\alpha)\otimes w$.  Pulling the lift through Mumford's
period charts, the cube isomorphism for translation gives
\[
             s'_{\alpha+y}=u_\alpha\,t_b^*s_\alpha,
 \qquad u_\alpha(0)=z.s(\alpha)\otimes w.
\]
The equality $[k]z=\iota(y)$, including its torus coordinates,
makes these formulas compatible with every period transition.
The zero-fibre ratio has exponents
$(k-1,1-k,1,-1)$ at $(0,b,c(\alpha),b+c(\alpha))$:
these are exactly the sum of the inverse-biextension exponents
$(-1,1,1,-1)$ and the fibre-line exponents $(k,-k,0,0)$.
Thus the formula is the source's \texttt{PhiRel}, with its
\emph{positive} index shift $+y$ and its exact zero normalization.
Any ratio-line generator with that zero value is the same generator,
since the abelian variety has only constant global units.  This
verifies the source's quantifier over all such generators, for
every theta lift and every section
\cite[Section 5.2.5]{dvr-unif-lan}.

This proves that the quotient point satisfies \texttt{QkToAkU}.
Conversely, if another $k$-torsion point satisfies that relation, its
projective theta action is the same, by the expansion bijection.
The difference has scalar theta action.  Commuting it with every
theta lift gives trivial Weil pairing with every $k$-torsion point.
The Weil pairing for the principal polarization in characteristic
zero is perfect, so the difference is zero.  This proves the exact
``if and only if'' in \texttt{TorsionIso}, also for $k=1$.
Comparing two orders of these same coefficient actions gives its
Weil-pairing field.  The convention is the original commutator of
lifts $t_x^*\mathscr L_K^k\to\mathscr L_K^k$, which equals the
polarization Weil pairing
\cite[Proposition 5.2.4.9]{dvr-unif-lan}.
Thus no pairing inversion or root-of-unity factor is omitted.

\subsection*{The exact integral theta module}

Since $\mathscr M$ is principal, $H^0(\mathcal B,\mathscr M)$ is
free of rank one over $R$, and a generator $\theta_0$ is nonzero on
the closed fibre.  At level one the period recurrence determines all
coefficients from coefficient zero.  By the expansion bijection,
there is a unique generic theta section $\theta$ with coefficient
zero $\theta_0$.  Its coefficient at $y$ is
$\psi_y\,\mathrm{rig}\,t_{c(y)}^*\theta_0$ in the original line
identifications.  Every coefficient is integral, and every nonzero
coefficient index is divisible by $\pi$.

We now prove regularity on the actual $\mathcal G$ without using
cohomology base change for this nonproper family.  Choose $m\geq0$
so that $\pi^m\theta$ extends to a section $\sigma$ on
$\mathcal G$.  Its coefficients are all divisible by $\pi^m$.
The exact expansion on the thickening $R/\pi^m$ says that $\sigma$
vanishes there.  Since $\mathcal G$ is flat over $R$ and
$\mathscr L$ is invertible, the sequence
\[
 0\longrightarrow\mathscr L\xrightarrow{\pi^m}\mathscr L
        \longrightarrow\mathscr L/\pi^m\mathscr L\longrightarrow0
\]
is exact.  Taking its kernel on global sections gives
$\sigma=\pi^m\widetilde\theta$ with
$\widetilde\theta\in H^0(\mathcal G,\mathscr L)$ and generic
restriction exactly $\theta$.  No scalar or extra power remains.

Every other integral section is $a\theta$ for some $a\in K$,
because the generic section space has dimension one.  Its coefficient
zero is $a\theta_0$, which is integral only if $a\in R$.
Thus $H^0(\mathcal G,\mathscr L)=R\theta$.  On the closed fibre all
coefficients except zero vanish, so $\theta$ restricts to
$q_B^*\theta_0$.  The map $q_B$ is surjective and $\theta_0$ is
nonzero, so the restriction is nonzero.  The generic restriction is
nonzero too.  These are precisely generation and fibrewise
nonvanishing in \texttt{PushInvertible} over $\operatorname{Spec}R$.

The theorem supplies the previously missing PEL degeneration and
torsion premises for the original complete-DVR arithmetic estimates.
Smooth formal coordinates at the identity supply their ordinary
formal truncations; compatibility with a chosen integral framed
chart uses its prescribed integral tangent basis.  The remaining
geometric construction gates are the coordinated finite level and
cover comparisons, global domination, and the height, section
remainder and cusp estimates.  No whole proof-map
node is declared closed by this local supplier.
